% !TEX program = xelatex
% !TEX encoding = UTF-8 Unicode
%
% Academic resume for Zhijun Yang.
% Layout adapted from Jake's Resume:
% https://github.com/jakegut/resume

\documentclass[a4paper,11pt]{article}

\usepackage{fontspec}
\setmainfont[
  BoldFont=texgyretermes-bold.otf,
  ItalicFont=texgyretermes-italic.otf,
  BoldItalicFont=texgyretermes-bolditalic.otf,
  Ligatures=NoCommon
]{texgyretermes-regular.otf}
\usepackage[
  a4paper,
  left=0.70in,
  right=0.70in,
  top=0.58in,
  bottom=0.58in
]{geometry}
\usepackage{enumitem}
\usepackage{tabularx}
\usepackage{titlesec}
\usepackage[hidelinks]{hyperref}
\usepackage{xcolor}

\hypersetup{
  pdftitle={Zhijun Yang -- Academic Resume},
  pdfauthor={Zhijun Yang},
  pdfsubject={Academic resume},
  pdfkeywords={distributed systems, disaggregated memory, CXL, RDMA, resource pooling}
}

\pagestyle{empty}
\setlength{\parindent}{0pt}
\setlength{\parskip}{0pt}
\setlength{\tabcolsep}{0pt}
\urlstyle{same}
\raggedright

\titleformat{\section}
  {\large\bfseries}
  {}
  {0pt}
  {}
  [\vspace{0.08em}\titlerule]
\titlespacing*{\section}{0pt}{0.78em}{0.42em}

\newcommand{\educationentry}[4]{%
  \begin{tabularx}{\textwidth}{@{}Xr@{}}
    \textbf{#1} & #2 \\
    \textit{#3} & \textit{#4}
  \end{tabularx}%
}

\newcommand{\datedentry}[2]{%
  \begin{tabular*}{\textwidth}{@{\extracolsep{\fill}}lr}
    #1 & #2
  \end{tabular*}%
}

\newcommand{\resumelink}[2]{\href{#1}{[#2]}}

\begin{document}

\begin{center}
  {\Huge\bfseries Zhijun Yang}\\[6pt]
  {\normalsize
    Ph.D. Student in Computer System Architecture,
    Huazhong University of Science and Technology
  }\\[3pt]
  {\small
    \href{mailto:cszjyang@hust.edu.cn}{cszjyang@hust.edu.cn}
    \enspace$\vert$\enspace
    \href{https://cszjyang.github.io/}{cszjyang.github.io}
    \enspace$\vert$\enspace
    \href{https://github.com/cszjyang}{github.com/cszjyang}
  }
\end{center}
\vspace{-0.55em}

\section{Education}

\educationentry
  {Huazhong University of Science and Technology (HUST)}
  {Wuhan, China}
  {Ph.D. in Computer System Architecture}
  {Sep. 2023 -- Jun. 2028 (expected)}

\vspace{0.15em}
\datedentry
  {\textit{Advisor:} \href{https://csyhua.github.io/csyhua/index.html}{Prof. Yu Hua}}
  {}

\vspace{0.30em}
\educationentry
  {Huazhong University of Science and Technology (HUST)}
  {Wuhan, China}
  {B.Eng. in Computer Science and Technology}
  {Sep. 2019 -- Jun. 2023}

\section{Research Interests}

Distributed systems, disaggregated memory, CXL/RDMA-based systems, and resource disaggregation and pooling.

\section{Publications}

\begin{enumerate}[
  leftmargin=1.55em,
  label={[\arabic*]},
  itemsep=0.42em,
  topsep=0.05em,
  parsep=0pt,
  partopsep=0pt
]

\item
  \textbf{Zhijun Yang}, Yu Hua, Ming Zhang, Menglei Chen, Xumin Chen, and Aoyang Tong.
  ``XTRA: Unifying Cache Coherence and Concurrency Control for Distributed Transactions in a CXL Pod.''
  \textit{32nd ACM Symposium on Operating Systems Principles (SOSP '26),} to appear.
  \resumelink{https://github.com/cszjyang/xtra}{Code}

\item
  \textbf{Zhijun Yang}, Yu Hua, Ming Zhang, Menglei Chen, and Yixiao Wang.
  ``FORGE: Mitigating Synchronization Amplification for Memory-Disaggregated Caching Systems.''
  \textit{20th USENIX Symposium on Operating Systems Design and Implementation (OSDI '26),}
  pp.~977--995, 2026.
  \textbf{Nominated for the Jay Lepreau Best Paper Award.}
  \resumelink{https://www.usenix.org/conference/osdi26/presentation/yang-zhijun}{Paper}
  \resumelink{https://github.com/cszjyang/forge}{Code}

\item
  Menglei Chen, Yu Hua, \textbf{Zhijun Yang}, Yixiao Wang, and Xumin Chen.
  ``CloverRec: Cost-Efficient Disaggregated Processing-in-Memory Systems for Deep Recommendation Inference.''
  \textit{25th USENIX Conference on File and Storage Technologies (FAST '27),} to appear.
  \resumelink{https://github.com/LighT-chenml/CloverRec}{Code}

\item
  Ming Zhang, Yu Hua, and \textbf{Zhijun Yang}.
  ``Motor: Enabling Multi-Versioning for Distributed Transactions on Disaggregated Memory.''
  \textit{18th USENIX Symposium on Operating Systems Design and Implementation (OSDI '24),}
  pp.~801--819, 2024.
  \resumelink{https://www.usenix.org/conference/osdi24/presentation/zhang-ming}{Paper}
  \resumelink{https://github.com/minghust/motor}{Code}

\end{enumerate}

\section{Academic Service}

\datedentry
  {Shadow Program Committee Member, EuroSys 2026 (Spring and Fall)}
  {2026}
\vspace{0.24em}
\datedentry
  {Artifact Evaluation Committee Member, USENIX ATC 2025}
  {2025}

\section{Honors and Awards}

\datedentry
  {Outstanding Graduate Student, HUST}
  {2024}
\vspace{0.20em}
\datedentry
  {Science and Technology Innovation Scholarship, HUST}
  {2024}
\vspace{0.20em}
\datedentry
  {Outstanding Graduate (B.Eng.), HUST}
  {2023}
\vspace{0.20em}
\datedentry
  {Award for Excellent Performance, NUS School of Computing Summer Workshop}
  {2021}

\end{document}
